\section{Forecasting with Flow Metrics}
Forecasting improves when you use cycle time and WIP. Flow metrics show whether the system can absorb new work without breaking.

\subsection{Jira setup steps (forecasting)}
\begin{enumerate}
  \item Enable Control Chart and Cumulative Flow Diagram on the board.
  \item Define cycle time boundaries (In Progress to Done).
  \item Set a target WIP level and monitor WIP breaches.
  \item Build a dashboard widget for average cycle time and throughput.
\end{enumerate}

Sample JQL filters:
\begin{itemize}
  \item Done last 30 days: \texttt{project = PMH AND status = Done AND resolved >= -30d}
  \item WIP breach candidates: \texttt{project = PMH AND status in (\"In Progress\", \"In Review\")}
\end{itemize}

\subsection{Worked example: forecasting a release window}
An Epic contains 24 Stories. The team last six sprints show 8–10 Stories finished per sprint. The PM forecasts a 3-sprint window with a 70\% confidence range, communicated as "3–4 sprints." Jira dashboards show cycle time and WIP; if WIP rises, the forecast widens.

\subsection{Case study insert}
Forecasting uses the completion rate of EPIC-4 and EPIC-5 stories. If cycle time rises, the R1 window is widened and the dashboard records the new range.
\subsection{Release burndown pitfalls revisited}
Use burndowns only for short horizons and stable scope. If scope shifts weekly, switch to throughput and cycle time to avoid false precision.










